\documentclass[11pt,letterpaper]{article}
\usepackage{../assets/scientific_report}
\usepackage[brazil]{babel}
\usepackage{amsmath,float,listings,xcolor,geometry,pgfplots}
\pgfplotsset{compat=1.18}
\geometry{headheight=23pt,top=2cm,bottom=2cm,left=2.5cm,right=2.5cm}
\setlength{\parskip}{0.5em}
\setlength{\parindent}{0pt}
\lstset{language=C,basicstyle=\ttfamily\footnotesize,keywordstyle=\color{blue}\bfseries,commentstyle=\color{green!60!black},numbers=left,numberstyle=\tiny\color{gray},numbersep=8pt,frame=single,breaklines=true,showstringspaces=false}

\begin{document}
\pagestyle{plain}
\begin{center}
{\LARGE\bfseries Programação em memória distribuída}\\[0.4em]
{\normalsize Eugenio Vitor Lopes dos Santos}\\[0.3em]
{\small DCA3703, Programação Paralela, Universidade Federal do Rio Grande do Norte}
\end{center}

\section{Enunciado}
Implementar uma troca de mensagens com exatamente dois processos MPI: o processo 0 envia uma mensagem ao processo 1, que responde com a mesma mensagem. Medir com \texttt{MPI\_Wtime} o tempo de múltiplas trocas consecutivas para tamanhos de 8 bytes a pelo menos 1 MB. Representar graficamente o tempo em função do tamanho e identificar as regiões dominadas pela latência e pelo custo de transferir os bytes.

\section{Fundamentação}
No MPI, cada processo possui seu próprio espaço de endereçamento. As rotinas relevantes para este experimento são:

\subsection{Inicialização e finalização}
\begin{itemize}
\item \texttt{MPI\_Init}: inicializa o ambiente MPI antes das demais chamadas.
\item \texttt{MPI\_Finalize}: encerra o ambiente MPI após a comunicação.
\end{itemize}

\subsection{Identificação dos processos}
\begin{itemize}
\item \texttt{MPI\_Comm\_rank}: informa o rank do processo em \texttt{MPI\_COMM\_WORLD}; aqui, 0 ou 1.
\item \texttt{MPI\_Comm\_size}: informa quantos processos há nesse comunicador; aqui, exatamente 2.
\end{itemize}

\subsection{Comunicação ponto a ponto}
\begin{itemize}
\item \texttt{MPI\_Send}: envia uma mensagem no modo padrão.
\item \texttt{MPI\_Recv}: recebe a mensagem correspondente ao envio.
\end{itemize}
No ping-pong, o rank 0 envia e depois recebe; o rank 1 recebe e responde com o mesmo buffer. Essa ordem evita deadlock.

\subsection{Outros modos de envio}
\begin{itemize}
\item \texttt{MPI\_Bsend}: envia usando um buffer fornecido pela aplicação com \texttt{MPI\_Buffer\_attach}.
\item \texttt{MPI\_Ssend}: envio síncrono; só conclui depois que o recebimento correspondente começa.
\end{itemize}
O programa não usa esses dois modos.

\section{Experimento}
O job 2109294 foi executado na partição \texttt{intel-512} do NPAD, em um nó identificado no log como \texttt{r1i3n15}. O Slurm reservou duas tarefas, uma CPU por tarefa. O programa foi compilado no job com \texttt{mpicc -std=c11 -O3 -Wall -Wextra} e iniciado com \texttt{mpirun -np 2}. O estado final foi \texttt{COMPLETED}, código \texttt{0:0}, em cinco segundos.

Para cada tamanho, ambos os processos alocam e preenchem um buffer. Uma \texttt{MPI\_Barrier} os sincroniza antes da medida. O processo 0 mede, com \texttt{MPI\_Wtime}, 10\,000 pares consecutivos de \texttt{MPI\_Send}/\texttt{MPI\_Recv}. Alocação, preenchimento, barreira, impressão e finalização do MPI ficam fora do intervalo. O tempo médio é o tempo total dividido por 10\,000 e corresponde à ida e volta completa, não a uma comunicação em apenas uma direção. Houve uma execução por tamanho; não foram medidas repetições independentes nem dispersão.

\section{Resultados}
A Tabela~\ref{tab:tempos} reproduz os valores do arquivo \texttt{tarefa14-2109294.out}. A Figura~\ref{fig:tempo} representa o tempo médio por troca; como o número de trocas é constante, o tempo total tem a mesma tendência.

\begin{table}[H]
\centering
\begin{tabular}{rrr}
\toprule
Mensagem (bytes) & Tempo total (s) & Tempo médio por troca ($\mu$s) \\
\midrule
8       & 0,006102265 & 0,610 \\
64      & 0,006119762 & 0,612 \\
512     & 0,008196224 & 0,820 \\
4\,096  & 0,048486192 & 4,849 \\
32\,768 & 0,120328724 & 12,033 \\
262\,144 & 0,811817241 & 81,182 \\
1\,048\,576 & 3,030861879 & 303,086 \\
\bottomrule
\end{tabular}
\caption{Medições do job 2109294, com 10\,000 trocas por tamanho.}
\label{tab:tempos}
\end{table}

\begin{figure}[H]
\centering
\begin{tikzpicture}
\begin{loglogaxis}[
 width=0.92\textwidth,height=7cm,
 xlabel={Tamanho da mensagem (bytes)},
 ylabel={Tempo médio por troca ($\mu$s)},
 xtick={8,64,512,4096,32768,262144,1048576},
 xticklabels={8,64,512,4 KiB,32 KiB,256 KiB,1 MiB},
 x tick label style={rotate=35,anchor=east,font=\small},
 grid=major,
 ymin=0.4,ymax=500,
 xmin=6,xmax=1400000,
]
\addplot+[mark=*,thick,color=primaryblue] coordinates {
 (8,0.610) (64,0.612) (512,0.820) (4096,4.849)
 (32768,12.033) (262144,81.182) (1048576,303.086)
};
\end{loglogaxis}
\end{tikzpicture}
\caption{Tempo médio de uma ida e volta em função do tamanho. Ambos os eixos usam escala logarítmica para mostrar os sete pontos.}
\label{fig:tempo}
\end{figure}

Entre 8 e 64 bytes, o tempo médio praticamente não muda: é a região em que a latência domina nesta medição. Em 512 bytes já há aumento, mas os sete pontos não localizam precisamente a transição. De 32 KiB a 1 MiB, o tempo cresce com o tamanho da mensagem, indicando maior influência da transferência dos bytes. Como os dois processos foram alocados no mesmo nó, estes dados não demonstram a largura de banda de uma rede entre nós.

\appendix
\clearpage
\section*{Anexo A: Código completo}
\lstinputlisting{tarefa14.c}

\section*{Anexo B: Script Slurm}
\lstinputlisting[language=bash]{tarefa14.slurm}

\end{document}
